% TOP_PROBLEM: {"id": 3243, "title": "Earth-Moon Problem", "category_id": 3, "category": {"id": 3, "name": "graph_theory", "display_name": "Graph Theory", "description": "Problems involving graphs, networks, and their properties.", "slug": "graph-theory", "order_index": 3, "created_at": "2026-07-31T15:26:25.671Z"}, "status": "open", "problem_number": "OPG-46900", "set_id": 12}
% FIRST_POSTED: 2026-10-01
% ENTRY_KIND: statement_only
% UnsolvedMath ID 3243; source catalogue code OPG-46900
% !TeX program = lualatex
% Definition and sources only; this file is not an LLM solution attempt.
\documentclass[11pt,a4paper]{article}
\usepackage[margin=25mm]{geometry}
\usepackage{fontspec}
\setmainfont{lmroman10-regular.otf}[BoldFont=lmroman10-bold.otf,
 ItalicFont=lmroman10-italic.otf,BoldItalicFont=lmroman10-bolditalic.otf]
\usepackage{amsmath,amssymb,mathtools}
\usepackage{unicode-math}
\setmathfont{latinmodern-math.otf}
\AtBeginDocument{\let\setminus\smallsetminus}
\usepackage{xurl,hyperref}
\hypersetup{colorlinks=true,urlcolor=blue}
\setcounter{secnumdepth}{0}
\setlength{\parindent}{0pt}
\setlength{\parskip}{5pt}
\setlength{\emergencystretch}{3em}
\begin{document}
\section*{Earth-Moon Problem}\label{3243}
{\small UnsolvedMath ID 3243 · OPG-46900 · Graph Theory · OpenGarden}
\subsection{Definitions and mathematical statement}
Let $G=(V,E)$ be a finite simple undirected graph. Say that $G$ has thickness
at most two if there are subsets $E_1,E_2\subseteq E$ with
$E=E_1\cup E_2$ such that both spanning graphs $(V,E_1)$ and $(V,E_2)$ are
planar. Planar means drawable in the plane with distinct vertices and
edge curves intersecting only at common endpoints. The two drawings are
independent: their vertices need not occupy the same geometric positions.
One may require $E_1\cap E_2=\varnothing$ by removing repeated edges from
one layer.

The chromatic number $\chi(G)$ is the least number of colours assigned to
vertices so that each edge has differently coloured ends. The Earth--Moon
problem is to determine the integer
\[
 M_2=\max\{\chi(G):G\text{ is finite simple and has thickness at most }2\}.
\]
This maximum is finite. Indeed, for each subgraph on $n\ge3$ vertices the
two planar layers together have at most $6n-12$ edges, so that subgraph has
a vertex of degree at most eleven. Successively deleting such vertices
and colouring in reverse order gives a twelve-colour upper bound.

The map interpretation identifies one country with one vertex. Adjacency
on either of the two worlds imposes an edge constraint, and both regions
of a country receive the same colour. Meeting at just an isolated boundary
point imposes no constraint. The graph formulation above specifies the
target without relying on geometric conventions about country shapes.
\subsection{Short English statement}
Determine the largest chromatic number of a graph whose edges can be split between two planar graphs.
\subsection{Sources}
\begin{itemize}
\item[S1] ULAM.AI and the UnsolvedMath Contributors, supplied problems.json, record 3243 (OPG-46900), OpenGarden collection. Catalogue identity and source transcription; CC BY 4.0.
\url{https://huggingface.co/datasets/ulamai/UnsolvedMath}.
\item[S2] Open Problem Garden, Earth-Moon Problem. Original problem and discussion; the mathematical formulation below is expanded from the supplied source record.
\url{http://www.openproblemgarden.org/op/earth_moon_problem}.
\end{itemize}
\end{document}
